% !PRINTED XII
\weberpage{XII}

\inhaltpar{87}{Erstes Beispiel: Kugelfunctionen}{273}
\inhaltpar{88}{Zweites Beispiel}{276}
\inhaltpar{89}{Die \gesperrt{Sturm}'schen Functionen}{279}
\inhaltpar{90}{\gesperrt{Hermite}'s Lösung des \gesperrt{Sturm}'schen Problems}{280}
\inhaltpar{91}{Bestimmung der \gesperrt{Hermite}'schen Form $H$}{282}
\inhaltpar{92}{Die Determinante der \gesperrt{Hermite}'schen Form}{283}
\inhaltpar{93}{Grundzüge der Charakteristikentheorie}{285}
\inhaltpar{94}{Charakteristik eines Systems von drei Functionen}{287}
\inhaltpar{95}{Beziehung der Charakteristiken zu den Schnittpunkten}{290}
\inhaltpar{96}{Anwendung der Charakteristiken auf die Eingrenzung der complexen Wurzeln einer Gleichung}{292}
\inhaltpar{97}{Bestimmung der Charakteristik}{294}
\inhaltpar{98}{\gesperrt{Gauss}' erster Beweis des Fundamentalsatzes der Algebra}{295}

\inhaltabschnitt{Neunter Abschnitt.}{Abschätzung der Wurzeln.}

\inhaltpar{99}{Das \gesperrt{Budan-Fourier}'sche Theorem}{299}
\inhaltpar{100}{Die \gesperrt{Newton}'sche Regel}{304}
\inhaltpar{101}{Der \gesperrt{Cartesi}'sche Lehrsatz}{308}
\inhaltpar{102}{Das \gesperrt{Jacobi}'sche Kriterium}{311}
\inhaltpar{103}{\gesperrt{Klein}'s geometrische Vergleichung der verschiedenen Kriterien}{312}
\inhaltpar{104}{Bestimmung einer oberen Grenze für die Wurzeln}{316}
\inhaltpar{105}{Abschätzung der imaginären Wurzeln}{318}
\inhaltpar{106}{Das Theorem von \gesperrt{Rolle}}{319}
\inhaltpar{107}{Die Sätze von \gesperrt{Laguerre} für Gleichungen mit nur reellen Wurzeln}{322}

\inhaltabschnitt{Zehnter Abschnitt.}{Genäherte Berechnung der Wurzeln.}

\inhaltpar{108}{Interpolation. Regula falsi}{331}
\inhaltpar{109}{Die \gesperrt{Newton}'sche Näherungsmethode}{335}
\inhaltpar{110}{Die Näherungsmethode von \gesperrt{Daniel Bernoulli} und verwandte Methoden}{341}
\inhaltpar{111}{Die Näherungsmethode von \gesperrt{Gräffe}}{344}
\inhaltpar{112}{Trigonometrische Auflösung cubischer Gleichungen}{349}
\inhaltpar{113}{Die \gesperrt{Gauss}'sche Methode der Auflösung trinomischer Gleichungen}{352}
\inhaltpar{114}{Berechnung der imaginären Wurzeln einer trinomischen Gleichung}{355}

\inhaltabschnitt{Elfter Abschnitt.}{Kettenbrüche.}

\inhaltpar{115}{Verwandlung rationaler Brüche in Kettenbrüche}{358}
\inhaltpar{116}{Kettenbruchentwickelung irrationaler Zahlen}{361}
\inhaltpar{117}{Die Näherungsbrüche}{362}
\inhaltpar{118}{Lösung unbestimmter Gleichungen aus zwei Unbekannten}{365}
